
\section{Finite Elements}

\subsection{Lines}

\subsubsection{LL2}
The linear two-noded line finite element has shape functions 
with the general form $\phi(\xi) = A + B \xi$.
\begin{figure}[htb]
  \begin{center}
    %\begin{overpic}[width=0.5\textwidth, grid, tics=10]{fig/LL2}
    %\begin{overpic}[width=0.5\textwidth,natwidth=360,natheight=108]{fig/LL2.pdf}
    \begin{overpic}[width=0.5\textwidth]{fig/LL2.pdf}
      \put(85, 22){$\xi$}
      \put(44, 25){$\xi=0$}
      %
      \put(18.5, 10){\bf 1}
      \put(17,  3){\footnotesize(-1)}
      %
      \put(78.5, 10){\bf 2}
      \put(77,  3){\footnotesize(1)}
    \end{overpic}
  \end{center}
  \caption{Linear two-noded line finite element.}
  \label{fig:LL2} % label must come after caption
\end{figure}

\begin{align}
 \phi_1(\xi) & = \;-\frac{1}{2} 
                 \hspace{0.65cm}
		 \cdot 
		 \hspace{0.65cm}
		 \left( \xi - 1 \right) \\
 \phi_2(\xi) & = \hspace{0.45cm} \frac{1}{2}
                 \;\left( \xi + 1 \right)
		 \hspace{0.65cm}
		 \cdot
\end{align}

\begin{Hremark}
\label{remark:node_connotation}
Note that 
``$\;\;\cdot$\;\;'' 
as used here has a dual purpose.  First, it denotes multiplication, and
by this measure alone, is unnecessary to include.  However, the second
purpose is a connotation of location, from left-to-right on the $\xi$-axis, 
where the shape function must be unity.  Moreover, the quantities in
parentheses 
(~\hspace{0.4cm}~) 
connote where the shape function must go to zero.  
\end{Hremark}

\clearpage
\subsubsection{LL3}
The linear three-noded line finite element has shape functions 
with the general form $\phi(\xi) = A + B \xi + C \xi^2$.
\begin{figure}[htb]
  \begin{center}
    %\begin{overpic}[width=0.5\textwidth, grid, tics=10]{fig/LL3}
    %\begin{overpic}[width=0.5\textwidth,natwidth=360,natheight=108]{fig/LL3.pdf}
    \begin{overpic}[width=0.5\textwidth]{fig/LL3.pdf}
      \put(85, 22){$\xi$}
      \put(44, 25){$\xi=0$}
      %
      \put(18.5, 10){\bf 1}
      \put(17,  3){\footnotesize(-1)}
      %
      \put(48.5, 10){\bf 2}
      \put(47.25,  3){\footnotesize(0)}
      %
      \put(78.5, 10){\bf 3}
      \put(77,  3){\footnotesize(1)}
    \end{overpic}
  \end{center}
  \caption{Linear three-noded line finite element.}
  \label{fig:LL3} % label must come after caption
\end{figure}

\begin{align}
 \phi_1(\xi) & = \;\;\frac{1}{2} 
                 \hspace{0.65cm}
		 \cdot 
		 \hspace{0.8cm}
                 \left( \xi  \right)
		 \;\;\left( \xi - 1 \right) \\
 \phi_2(\xi) & = -1 \;
                 \left( \xi + 1 \right)
		 \hspace{0.42cm}
		 \cdot
		 \hspace{0.42cm} 
		 \left( \xi - 1 \right) \\
 \phi_3(\xi) & = \;\;\frac{1}{2} \;
                 \left( \xi + 1 \right)
		 \;\;\;\left( \xi  \right)
		 \hspace{0.7cm}
		 \cdot
\end{align}
See Remark~\ref{remark:node_connotation} for an explanation of the 
multiplication 
``$\;\;\cdot$\;\;'' 
and parentheses (~\hspace{0.4cm}~) notation.



\clearpage
\subsubsection{LL4}
The linear four-noded line finite element has shape functions 
with the general form $\phi(\xi) = A + B \xi + C \xi^2 + D \xi^3$.
\begin{figure}[htb]
  \begin{center}
    %\begin{overpic}[width=0.5\textwidth, grid, tics=10]{fig/LL4}
    %\begin{overpic}[width=0.5\textwidth,natwidth=360,natheight=108]{fig/LL4.pdf}
    \begin{overpic}[width=0.5\textwidth]{fig/LL4.pdf}
      \put(85, 22){$\xi$}
      \put(44, 25){$\xi=0$}
      %
      \put(18.5, 10){\bf 1}
      \put(17,  3){\footnotesize(-1)}
      %
      \put(38.5, 10){\bf 2}
      \put(35,  3){\footnotesize$\left(-\frac{1}{3}\right)$}
      %
      \put(58.5, 10){\bf 3}
      \put(56.5,  3){\footnotesize$\left(\frac{1}{3}\right)$}
      %
      \put(78.5, 10){\bf 4}
      \put(77,  3){\footnotesize(1)}
    \end{overpic}
  \end{center}
  \caption{Linear four-noded line finite element.}
  \label{fig:LL4} % label must come after caption
\end{figure}

\begin{align}
 \phi_1(\xi) & = \frac{-9}{16} 
                 \hspace{0.5cm}
		 \cdot 
		 \hspace{0.5cm}
                 \left( \xi + \frac{1}{3} \right)
                 \left( \xi - \frac{1}{3} \right) 
		 \left( \xi - 1 \right) \\
 \phi_2(\xi) & = \frac{9}{16} \;
                 \left( \xi + 1 \right)
		 \hspace{0.7cm}
		 \cdot
		 \hspace{0.7cm} 
		 \left( \xi - \frac{1}{3} \right) 
		 \left( \xi - 1 \right) \\
 \phi_3(\xi) & = \frac{-9}{16} 
                 \left( \xi + 1 \right)
		 \left( \xi + \frac{1}{3} \right)
		 \hspace{0.7cm}
		 \cdot
		 \hspace{0.7cm}
		 \left( \xi - 1 \right) \\
 \phi_4(\xi) & = \frac{9}{16} \;
                 \left( \xi + 1 \right)
		 \left( \xi + \frac{1}{3} \right)
		 \left( \xi - \frac{1}{3} \right)
		 \hspace{0.5cm}
		 \cdot
\end{align}
See Remark~\ref{remark:node_connotation} for an explanation of the 
multiplication 
``$\;\;\cdot$\;\;'' 
and parentheses (~\hspace{0.4cm}~) notation.


\clearpage
\subsection{Triangles}

\subsubsection{LT3}
The linear three-noded triangular finite element has shape functions 
with the general form $\phi(\vxi) = A + B \xi_1 + C \xi_2$.
\begin{figure}[htb]
  \begin{center}
    %\begin{overpic}[width=0.5\textwidth, grid, tics=10]{fig/LT3}
    %\begin{overpic}[width=0.5\textwidth,natwidth=360,natheight=360]{fig/LT3.pdf}
    \begin{overpic}[width=0.5\textwidth]{fig/LT3.pdf}
      \put(85, 20){$\xi_1$}
      \put(20, 85){$\xi_2$}
      %
      \put( 8,  8){\bf 3}
      \put( 4,  3){\footnotesize(0, 0)}
      %
      \put(80,  8){\bf 1}
      \put(76,  3){\footnotesize(1, 0)}
      %
      \put( 8, 80){\bf 2}
      \put( 4, 75){\footnotesize(0, 1)}
    \end{overpic}
  \end{center}
  \caption{Linear three-noded triangular finite element.}
  \label{fig:LT3} % label must come after caption
\end{figure}

\begin{align}
 \phi_1(\vxi) & = \xi_1 \\
 \phi_2(\vxi) & = \xi_2 \\
 \phi_3(\vxi) & = 1 - \xi_1 - \xi_2 
\end{align}

\clearpage
\subsubsection{LT6}
The linear six-noded triangular finite element has shape functions 
with the general form $\phi(\vxi) = A + B \xi_1 + C \xi_2 + D \xi_1 \xi_2
 + E \xi_1^2 \xi_2 + F \xi_1 \xi_2^2.$
\begin{figure}[htb]
  \begin{center}
    %\begin{overpic}[width=0.5\textwidth, grid, tics=10]{fig/LT6}
    %\begin{overpic}[width=0.5\textwidth,natwidth=360,natheight=360]{fig/LT6.pdf}
    \begin{overpic}[width=0.5\textwidth]{fig/LT6.pdf}
      \put(85, 20){$\xi_1$}
      \put(20, 85){$\xi_2$}
      %
      \put( 8,  8){\bf 3}
      \put( 4,  3){\footnotesize(0, 0)}
      %
      \put(80,  8){\bf 1}
      \put(76,  3){\footnotesize(1, 0)}
      %
      \put( 8, 80){\bf 2}
      \put( 4, 75){\footnotesize(0, 1)}
      %
      \put(55, 55){\bf 4}
      \put(50, 50){\footnotesize($\half$, $\half$)}
      %
      \put( 8, 45){\bf 5}
      \put( 4, 40){\footnotesize(0, $\half$)}
      %
      \put(46,  8){\bf 6}
      \put(42,  3){\footnotesize($\half$, 0)}
    \end{overpic}
  \end{center}
  \caption{Linear six-noded triangular finite element.}
  \label{fig:LT6} % label must come after caption
\end{figure}

\begin{align}
 \phi_1(\vxi) & = 2 \; \xi_1 \left(\xi_1 - \half \right) \\
 \phi_2(\vxi) & = 2 \; \xi_2 \left(\xi_2 - \half \right) \\
 \phi_3(\vxi) & = 2 \left( 1 - \xi_1 - \xi_2 \right) 
                  \left( \half - \xi_1 - \xi_2 \right) \\
 \phi_4(\vxi) & = 4 \; \xi_1 \; \xi_2 \\
 \phi_5(\vxi) & = 4 \left(1 - \xi_1 - \xi_2\right) \xi_2 \\
 \phi_6(\vxi) & = 4 \left(1 - \xi_1 - \xi_2\right) \xi_1
\end{align}

\clearpage
\subsubsection{LT6 Detailed Derivation}

\begin{figure}[htb]
  \begin{center}
    %\begin{overpic}[width=\textwidth, grid, tics=10]{fig/LT6_explicated}
    %\begin{overpic}[width=\textwidth,natwidth=684,natheight=360]{fig/LT6_explicated.pdf}
    \begin{overpic}[width=\textwidth]{fig/LT6_explicated.pdf}
      \put(50, 10){$\xi_1$}
      \put(96, 10){$\xi_1$}
      \put(10, 49){$\xi_2$}
      \put(57, 49){$\xi_2$}
      %
      \put(9.5,  2){$H_1$}
      \put(60,  2){$H_2$}
      %
      \put( 2, 32.5){$V_1$}
      \put(82, 26){$V_2$}
      %
      \put(25.3, 44.8){$D_1$}
      \put(70, 41){$D_2$}
    \end{overpic}
  \end{center}
  \caption{Detailed derivation of the linear 
  six-noded triangular finite element, with horizontal lines $H_1$ and 
  $H_2$, vertical lines $V_1$ and $V_2$, and diagonal lines $D_1$ and 
  $D_2$.}
  \label{fig:LT6_explicated} % label must come after caption
\end{figure}

We begin by writing equations of two horizontal, two vertical, and
two diagonal lines, as outlined in Fig.~(\ref{fig:LT6_explicated}).
\begin{align}
 H_1 & : \xi_2 = 0 \;\;\;\forall\;\;\; \xi_1 \\
 H_2 & : \xi_2 = \half \;\;\forall\;\;\; \xi_1  
         \implies 0 = \left( \xi_2 - \half \right) \\
 V_1 & : \xi_1 = 0 \;\;\;\forall\;\;\; \xi_2 \\
 V_2 & : \xi_1 = \half \;\;\forall\;\;\; \xi_2 
         \implies 0 = \left( \xi_1 - \half \right) \\
 D_1 & : \xi_2 = 1 - \xi_1 \;\;\;
         \implies 0 = \left( 1 - \xi_1 - \xi_2 \right) \\
 D_2 & : \xi_2 = \half - \xi_1 \;\;
         \implies 0 = \left( \half - \xi_1 - \xi_2 \right)
\end{align}
These six equations represent lines where the shape function has zero
value for all combinations of $(\xi_1, \xi_2)$.  Using these lines, 
we can construct shape functions as a functional products of $(\xi_1, \xi_2$)
and locations (e.g., lines that intersect nodes) where the nodal
values of the shape function must be identically zero.  

Let shape function $\phi_1$ have coefficient $A_1$.  Then $\phi_1$ can be
written as 
\be
\phi_1(\vxi) = A_1 \cdot V_1 \cdot V_2 = A_1 \cdot \xi_1 \cdot
\left(\xi_1 - \half \right).
\ee
Then, evaluating $\phi_1$ at the node {\bf 1} location gives
\be
\phi_1(1,0) = A_1 (1) \left( \half \right) \overset{\mbox{\tiny set}}{=} 1
\implies A_1 = 2.
\ee
The form of $\phi_2$ follows similarly to $\phi_1$.
Next, consider $\phi_3$.
\be 
\phi_3(\vxi) = A_3 \cdot D_1 \cdot D_2 = A_3 \cdot \left( \half 
- \xi_1 - \xi_2 \right) \cdot \left( 1 - \xi_1 - \xi_2 \right).
\ee
Then, evaluating $\phi_3$ at the node {\bf 3} location gives
\be
\phi_3(0,0) = A_3 \left( \half \right) (1) \overset{\mbox{\tiny set}}{=} 1
\implies A_3 = 2.
\ee
Finally, $\phi_4$, $\phi_5$ and $\phi_6$ can be found readily (we omit
some details now that the pattern has been shown above)
as
\begin{align}
\phi_4 & = A_4 \cdot H_1 \cdot V_1 \;\;\;\;\mbox{and}\;\;\; A_4 = 4, \\
\phi_5 & = A_5 \cdot H_1 \cdot D_1 \;\;\;\mbox{and}\;\;\; A_5 = 4, \\
\phi_6 & = A_6 \cdot V_1 \cdot D_1 \;\;\;\;\mbox{and}\;\;\; A_6 = 4. 
\end{align}



\clearpage
\subsubsection{LT10}
The linear ten-noded triangular finite element has shape functions 
with the general form $\phi(\vxi) = A + B \xi_1 + C \xi_2 + D \xi_1 \xi_2
 + E \xi_1^2 + F \xi_2^2 + G \xi_1^2 \xi_2 + H \xi_1 \xi_2^2 + I \xi_1^3
 + J \xi_2^3$.
\begin{figure}[htb]
  \begin{center}
    %\begin{overpic}[width=0.5\textwidth, grid, tics=10]{fig/LT10}
    %\begin{overpic}[width=0.5\textwidth,natwidth=360,natheight=360]{fig/LT10.pdf}
    \begin{overpic}[width=0.5\textwidth]{fig/LT10.pdf}
      \put(85, 20){$\xi_1$}
      \put(20, 85){$\xi_2$}
      %
      \put( 8,  8){\bf 3}
      \put( 4,  3){\footnotesize(0, 0)}
      %
      \put(80,  8){\bf 1}
      \put(76,  3){\footnotesize(1, 0)}
      %
      \put( 8, 80){\bf 2}
      \put( 4, 75){\footnotesize(0, 1)}
      %
      \put( 8, 55){\bf 6}
      \put( 4, 50){\footnotesize(0, $\frac{2}{3}$)}
      %
      \put( 8, 35){\bf 7}
      \put( 4, 30){\footnotesize(0, $\frac{1}{3}$)}
      %
      \put(38, 38){\bf 10}
      \put(38, 29){\footnotesize($\frac{1}{3}$, $\frac{1}{3}$)}
      % 
      \put(58, 38){\bf 4}
      \put(64, 38){\footnotesize($\frac{2}{3}$, $\frac{1}{3}$)}
      %
      \put(58,  8){\bf 9}
      \put(54,  3){\footnotesize($\frac{2}{3}$, 0)}
      %
      \put(38, 58){\bf 5}
      \put(44, 58){\footnotesize($\frac{1}{3}$, $\frac{2}{3}$)}
      %
      \put(38,  8){\bf 8}
      \put(34,  3){\footnotesize($\frac{1}{3}$, 0)}
      %
    \end{overpic}
  \end{center}
  \caption{Linear ten-noded triangular finite element.}
  \label{fig:LT10} % label must come after caption
\end{figure}

\begin{align}
 \phi_1(\vxi) & = \frac{9}{2} \; \xi_1 
                  \left(\xi_1 - \frac{1}{3} \right)
		  \left(\xi_1 - \frac{2}{3} \right) \\
 \phi_2(\vxi) & = \frac{9}{2} \; \xi_2 
                  \left(\xi_2 - \frac{1}{3} \right) 
		  \left(\xi_2 - \frac{2}{3} \right) \\
 \phi_3(\vxi) & = \frac{9}{2} 
                  \left( 1 - \xi_1 - \xi_2 \right) 
                  \left( \frac{1}{3} - \xi_1 - \xi_2 \right) 
		  \left( \frac{2}{3} - \xi_1 - \xi_2 \right) \\
 \phi_4(\vxi) & = \frac{27}{2} \; \xi_1 \; \xi_2 
                  \left( \xi_1 - \frac{1}{3} \right) \\
 \phi_5(\vxi) & = \frac{27}{2} \; \xi_1 \; \xi_2 
                  \left( \xi_2 - \frac{1}{3} \right) \\
 \phi_6(\vxi) & = \frac{27}{2} 
                  \left(1 - \xi_1 - \xi_2\right) 
		  \left( \xi_2 - \frac{1}{3} \right) \xi_2 \\
 \phi_7(\vxi) & = \frac{27}{2}
                  \left( 1 - \xi_1 - \xi_2 \right)
		  \left( \frac{2}{3} - \xi_1 - \xi_2 \right) \; \xi_2 \\
 \phi_8(\vxi) & = \frac{27}{2} 
                  \left( 1 - \xi_1 - \xi_2 \right) 
		  \left( \frac{2}{3} - \xi_1 - \xi_2 \right) \; \xi_1 \\
 \phi_9(\vxi) & = \frac{27}{2}
                  \left(1 - \xi_1 - \xi_2\right)
		  \left( \xi_1 - \frac{1}{3} \right) \xi_1 \\
 \phi_{10}(\vxi) & = 27 \left( 1 - \xi_1 - \xi_2 \right ) \xi_1 \; \xi_2
\end{align}


\clearpage
\subsection{Quadrilaterals}

\subsubsection{LQ4}
The linear four-noded quadrilateral finite element has shape functions 
with the general form $\phi(\vxi) = A + B \xi_1 + C \xi_2 + D \xi_1 \xi_2$.
\begin{figure}[htb]
  \begin{center}
    %\begin{overpic}[width=0.4\textwidth, grid, tics=10]{fig/LQ4}
    %\begin{overpic}[width=0.4\textwidth,natwidth=288,natheight=288]{fig/LQ4.pdf}
    \begin{overpic}[width=0.4\textwidth]{fig/LQ4.pdf}
      \put(85, 55){$\xi_1$}
      \put(55, 85){$\xi_2$}
      %
      \put(18, 18){\bf 1}
      \put(12, 10){\footnotesize(-1, -1)}
      %
      \put(82, 18){\bf 2}
      \put(78, 10){\footnotesize(1, -1)}
      %
      \put(82, 80){\bf 3}
      \put(78, 88){\footnotesize(1, 1)}
      %
      \put(18, 80){\bf 4}
      \put(12, 88){\footnotesize(-1, 1)}
    \end{overpic}
  \end{center}
  \caption{Linear four-noded quadrilateral finite element.}
  \label{fig:LQ4}  % label must come after caption
\end{figure}

\begin{align}
 \phi_1(\vxi) & = \frac{1}{4} \left( 1 - \xi_1 \right) \left( 1 - \xi_2 \right) 
 \\
 \phi_2(\vxi) & = \frac{1}{4} \left( 1 + \xi_1 \right) \left( 1 - \xi_2 \right) 
 \\
 \phi_3(\vxi) & = \frac{1}{4} \left( 1 + \xi_1 \right) \left( 1 + \xi_2 \right) 
 \\
 \phi_4(\vxi) & = \frac{1}{4} \left( 1 - \xi_1 \right) \left( 1 + \xi_2 \right) 
\end{align}




\clearpage
\subsubsection{LQ8}
The quadratic eight-noded quadrilateral finite element has shape functions
with the general form 
$\phi(\vxi) = A + B \xi_1 + C \xi_2 + D \xi_1 \xi_2 + E \xi_1^2 +
F \xi_2^2 + G \xi_1^2 \xi_2 + H \xi_1 \xi_2^2$.
This is often referred to as the ``serendipity'' quadrilateral element
[explanation to come].

\begin{figure}[htb]
  \begin{center}
    %\begin{overpic}[width=0.4\textwidth, grid, tics=10]{fig/LQ8}
    %\begin{overpic}[width=0.4\textwidth,natwidth=288,natheight=288]{fig/LQ8.pdf}
    \begin{overpic}[width=0.4\textwidth]{fig/LQ8.pdf}
      \put(80, 55){$\xi_1$}
      \put(55, 80){$\xi_2$}
      %
      \put(18, 18){\bf 1}
      \put(12, 10){\footnotesize(-1, -1)}
      %
      \put(82, 18){\bf 2}
      \put(78, 10){\footnotesize(1, -1)}
      %
      \put(82, 80){\bf 3}
      \put(78, 88){\footnotesize(1, 1)}
      %
      \put(18, 80){\bf 4}
      \put(12, 88){\footnotesize(-1, 1)}
      %
      \put(48, 15){\bf 5}
      \put(43,  7){\footnotesize(0, -1)}
      %
      \put(92, 50){\bf 6}
      \put(88, 42){\footnotesize(1, 0)}
      %
      \put(48, 90){\bf 7}
      \put(43, 98){\footnotesize(0, 1)}
      %
      \put(13, 50){\bf 8}
      \put( 6, 42){\footnotesize(-1, 0)}
    \end{overpic}
  \end{center}
  \caption{Linear eight-noded quadrilateral finite element.}
  \label{fig:LQ8} % label must come after caption
\end{figure}


\begin{align}
 \phi_1(\vxi) & = \frac{1}{4} \left( 1 - \xi_1 \right) \left( 1 - \xi_2 \right) 
 \left(-1 - \xi_1 - \xi_2 \right)
 \\
 \phi_2(\vxi) & = \frac{1}{4} \left( 1 + \xi_1 \right) \left( 1 - \xi_2 \right) 
 \left(-1 + \xi_1 - \xi_2 \right)
 \\
 \phi_3(\vxi) & = \frac{1}{4} \left( 1 + \xi_1 \right) \left( 1 + \xi_2 \right) 
 \left(-1 + \xi_1 + \xi_2 \right)
 \\
 \phi_4(\vxi) & = \frac{1}{4} \left( 1 - \xi_1 \right) \left( 1 + \xi_2 \right) 
 \left(-1 - \xi_1 + \xi_2 \right)
 \\
 \phi_5(\vxi) & = \half \left( 1 - \xi_1^2 \right) \left( 1 - \xi_2 \right)
 \\
 \phi_6(\vxi) & = \half \left( 1 + \xi_1 \right) \left( 1 - \xi_2^2 \right)
 \\
 \phi_7(\vxi) & = \half \left( 1 - \xi_1^2 \right) \left( 1 + \xi_2 \right)
 \\
 \phi_8(\vxi) & = \half \left( 1 - \xi_1 \right) \left( 1 - \xi_2^2 \right)
\end{align}
